\chapter{The Present Research}\label{chap:dolor}
Persistent gender disparities in mathematics continue to influence educational and career choices, often deterring girls from fully engaging with the subject. Prior research, as discussed in \ref{chap:SocialScience} has demonstrated that gender stereotypes and stereotype threat adversely affect both performance and self-efficacy in mathematics. The present study examines whether a brief, school-based intervention can diminish stereotype endorsement and promote greater confidence among female students.

\section{Research Questions}
The research questions are the following:
\begin{itemize}
    \item[RQ1.] \textbf{Is it possible to reduce gender stereotypes related to mathematics through the proposed intervention?} \\ This question aims to investigate whether exposure to female role models in mathematics can specifically influence students' beliefs about gender and mathematical ability within the school context.
    \item[RQ2.] \textbf{Is it possible to increase females' self-efficacy in mathematics within the context of the present study?} \\ Here the focus is on whether the intervention can enhance girls' confidence in their own mathematical skills. By providing relatable role models, students' belief in their capacity to succeed in math-related tasks might be strengthened.
    \item[RQ3.] \textbf{Is it possible to decrease females' math anxiety by means of the school-based intervention?} \\ This question examines whether exposure to positive examples of women succeeding in mathematics can reduce girls' apprehension toward math. Reducing anxiety is crucial, as high levels of math anxiety can negatively affect performance and engagement.
    \item[RQ4.] \textbf{Is it possible to improve students' attitudes toward mathematics following the intervention?} \\ The final question explores whether the intervention can foster more positive attitudes toward the subject, such as interest, enjoyment, and motivation. Seeing female role models may not only counter stereotypes but also make math appear more accessible and relevant to all students.
\end{itemize}

\section{Research Design}
This study employs a short-term quantitative experimental design using Likert-scale questionnaires. A within-group design is adopted, where the same students are assessed before and after a seminar featuring female mathematicians to evaluate changes in gender stereotypes, self-efficacy, math anxiety, and attitudes toward mathematics.

\section{Sample}
The research involved 362 students from a middle school (\textit{scuola secondaria di primo grado}) located in the city center of Trieste. The school serves a total population of approximately 533 students.

The school was selected using purposive sampling due to its specific characteristics and availability for the project. Similarly, the classes participating in the experiment were selected through purposive sampling, based on teacher willingness and scheduling compatibility.

The sample was distributed across three grade levels. Specifically, the participants included 127 students from Grade 6 (63 males, 62 females, and 2 who preferred not to disclose their gender), 121 students from Grade 7 (59 males and 62 females), and 114 students from Grade 8 (64 males and 50 females). Participants ranged in age from 10 to 14 years.

\section{Procedure}
The study employed a pre- and post-test within-group experimental design to evaluate the effects of exposure to female role models in mathematics. The data collection and intervention took place in December 2025. The experimental procedure was divided into two distinct sessions to assess the retention of the intervention's effects.

The first session was conducted in the school's auditorium (aula magna), where the researcher personally welcomed groups of three or four classes at a time. This session lasted approximately 60 minutes and began with the administration of the pre-test questionnaire (approx. 20 minutes). Immediately following the data collection, the researcher delivered the intervention.

The seminar opened with an interactive three-step introduction designed to stimulate critical thinking. First, students were asked to write down the names of famous mathematicians on a slip of paper. Second, the researcher displayed images of well-known mathematicians, asking the students to identify what these figures had in common. Finally, the class was engaged in a discussion regarding the reasons behind the apparent lack of female representation in the field.

Following this introductory discussion, the researcher presented a chronological overview of female contributions to mathematics. The narrative covered figures such as Hypatia of Alexandria, the pioneers of the 18th and 19th centuries (Maria Gaetana Agnesi, Sophie Germain, and Sofya Kovalevskaya), and the \textit{Human Computers} of the 20th century. Throughout this historical section, the presentation explicitly highlighted the social and institutional obstacles these women faced, such as exclusion from universities and the necessity of using male pseudonyms to be taken seriously.

Particular emphasis was placed on the modern era and the life of Maryam Mirzakhani. The presentation detailed her achievement as the first female Fields Medalist. To foster a sense of relatability, the researcher highlighted her personal journey, noting that she originally dreamt of becoming a writer rather than a mathematician and mentioning her participation in the International Mathematical Olympiad.Finally, her complex theoretical work on curved surfaces was explained with an age-appropriate manner and directly linked to tangible modern applications such as GPS technology and Artificial Intelligence to demonstrate the practical utility of advanced mathematics.

The second session took place exactly one week after the intervention. In this final phase, participants completed the post-test questionnaire—identical to the pre-test—to assess changes in their perceptions over time.

\section{Instruments}\label{sec:Instruments}
A questionnaire was used to collect data. 
The complete version of the questionnaire is provided in \autoref{app:questionnaire}.
The questionnaire consisted of four sections, each designed to measure a specific construct relevant to the study; all items were rated on a five-point Likert scale ranging from 1 to 5.
\begin{itemize}
    \item Math Anxiety.
    This section includes items assessing feelings of tension, nervousness, or discomfort that may arise when performing mathematical tasks. The items were adapted from short versions of existing math anxiety scales, focusing on students’ emotional reactions during tests or classroom activities.
    \item Mathematics Self-Efficacy. 
    This section measures students’ perceived confidence in their mathematical abilities—that is, their belief in being able to successfully complete exercises and understand new concepts. The items were formulated to assess self-efficacy beliefs applied specifically to mathematics learning.
    \item Attitudes Toward Mathematics.
    This section explores students’ general attitudes toward mathematics, including interest, enjoyment, and perceived usefulness of the subject. The items were inspired by the Fennema-Sherman Mathematics Attitudes Scales (\cite{FennemaSherman1976}) and by \textcite{Passolunghi2014}, which investigate motivational and affective factors linked to mathematics learning and gender differences.
    \item Gender Stereotypes in Mathematics.
    The final section investigates students’ beliefs and stereotypes regarding gender differences in mathematical ability. The items were developed based on and inspired by studies by \textcite{Bedynska2021} and \textcite{Dersch2022}, which analyze the influence of stereotype threat and math-gender misconceptions on students’ performance and self-perception.
\end{itemize}


\section{Data Collection}
Data were collected during regular school hours in the participating middle schools in the Province of Trieste. The pre-test and post-test questionnaires were administered in paper format by the researcher, with the assistance of classroom teachers. Each session lasted approximately 10 minutes.

The pre-test was conducted immediately before the seminar, and the post-test was completed right after the seminar ended. Students filled out the questionnaires individually and anonymously. Teachers and the researcher were present to provide clarification if needed, while ensuring that responses remained private and that students did not discuss their answers with peers.

All completed questionnaires were collected on site, coded for analysis, and securely stored in accordance with ethical and data protection standards.

\section{Data Analysis}
\section{Ethical Consideration}
This research was carried out in full compliance with recognized ethical principles for studies involving human participants, particularly minors. It followed the guidelines of the Declaration of Helsinki (1964), the European Code of Conduct for Research Integrity, and the Code of Ethics and Code of Conduct of the University of Trieste. Before the intervention and data collection began, formal authorization was obtained from the school’s Principal, and informed consent was secured from both students and their parents or legal guardians. Participation was entirely voluntary, and students were clearly informed that they could withdraw from the study at any point without any negative consequences for their academic evaluation or school experience.

To protect anonymity and confidentiality, no personally identifying data were gathered. All questionnaire responses were coded, stored securely, and analyzed exclusively in aggregate form. Neither the identities of the students nor those of their schools were revealed in any reports, publications, or presentations arising from this research.

Particular attention was devoted to ensuring that all materials and procedures were age-appropriate and presented in a manner that avoided discomfort or confusion. Since the study addressed issues related to gender and mathematics, care was taken to ensure that no participant felt singled out or judged based on gender. Teachers were informed in advance about the goals and structure of the intervention and were present during all sessions to provide support if necessary. The seminar was framed positively, emphasizing inclusivity, diversity, and the accomplishments of women in mathematics, with the broader aim of creating a respectful and encouraging classroom environment.